\section{What Remains at Quadratic Length}
\label{sec:quadratic-diagnostics}

Here $P=P_{N,m}$ is the half-lazy adjacent pair heat bath on the full
uniform identity fiber, $m=2\lfloor N^2/2\rfloor$, and $N$ is sufficiently
large.  Write $G_j=t_1\cdots t_j$.  The results below are proved with
explicit comparisons.  They give polynomial lower bounds on mixing time
and rule out two broad locations for a superpolynomial obstruction; they
do not decide whether the full quadratic-length walk mixes polynomially.

\begin{theorem}[Cuts observing well-separated prefixes]
\label{thm:spaced-prefix-cuts}
Put $B_0=\lceil4N\log N\rceil$.  Suppose
$0=j_0<j_1<\cdots<j_k<j_{k+1}=m$ and every gap is at least $B_0+2$.
For every nontrivial event $A$ determined by the entire tuple
$(G_{j_1},\ldots,G_{j_k})$,
\[
 \frac{Q_P(A,A^c)}{\min\{\pi(A),\pi(A^c)\}}
       \ge\frac{c}{mN}.                            \tag{QD.1}
\]
This includes arbitrarily rare cuts and $k=\Theta(N/\log N)$ observations.
The order is sharp: the midpoint derangement event has constant mass on
both sides and cut ratio $\Theta(1/(mN))=\Theta(N^{-3})$.
\end{theorem}
\begin{proof}
Equation (CB.7) bounds the relative parity-class error of every segment
of length at least $B_0$ by $\epsilon_N\le2/N^2$.  Consequently the
joint prefix law has density bounded above and below by absolute constants
relative to independent uniform permutations of the prescribed parities:
its density is the product of $k+1$ factors $1\pm\epsilon_N$, divided by
the full bridge normalization.  Here $k=O(N/\log N)$, so the product
does not accumulate an unbounded loss.

Fix an observed bond $j=j_r$, and expose its adjacent gates $a=t_j$ and
$b=t_{j+1}$.  Removing them shortens the two neighboring segments by one.
The same estimate shows that the joint law of the prefix tuple and $a,b$
is comparable, by absolute constants, to the preceding independent
prefix law times two independent uniform transpositions.  This comparison
is uniform over the tuple and the exposed gates, even for rare events.

The Hurwitz replacement $(a,b)\mapsto(b,bab)$ changes only the observed
prefix $G_j$, to $G_jab$.  When $a\ne b$ its conditional heat-bath
probability is at least $1/3$; when $a=b$ the proposed displacement is
zero.  Thus the Dirichlet form of $P$ on functions of the observed tuple
dominates $c/(m-1)$ times the sum of the forms that multiply one
independent parity-uniform coordinate by two uniform transpositions.
Each such two-step permutation walk has gap at least $c/N$: the
contents-ratio bound used in (CB.7), after pairing conjugate diagrams,
bounds every nontrivial absolute eigenvalue of the one-step walk by
$1-1/(N-1)$.  Tensorization and the bounded density comparison give
\[
 \mathcal E_P(F,F)\ge\frac{c}{mN}\operatorname{Var}_\pi F.
\]
For an indicator, variance is at least half the smaller-side mass,
which proves (QD.1).

At the midpoint the parity-uniform derangement probability tends to
$e^{-1}$, and the bridge density tends uniformly to one.  Only the
midpoint bond can cross this cut.  Given its selection, a replacement
changes the prefix on at most four inputs.  Under the comparable law in
which the prefix is parity-uniform and the adjacent gates are independent,
the probability that any affected input is fixed before or after the
change is $O(1/N)$.  The same upper bound holds under the bridge law.
This gives the matching upper bound; (QD.1) gives the lower bound.
\end{proof}

Identify words under simultaneous conjugation by $S_N$, retaining their
gate order and incidence pattern.  The pair kernel is equivariant, so it
defines a genuine quotient chain $\bar P$.  Its stationary orbit weights
are orbit cardinalities divided by $|\Omega(N,m)|$, not uniform weights
on isomorphism classes.

\begin{theorem}[Reduction to relabeling-invariant cuts]
\label{thm:relabeling-reduction}
At the above quadratic length,
\[
 c\min\{\Phi(\bar P),N^{-11}\}
       \le\Phi(P)\le\Phi(\bar P).                 \tag{QD.2}
\]
Thus superpolynomially small full conductance is equivalent to
superpolynomially small conductance of a relabeling-invariant cut.
The exponent $11$ is a conservative path-counting bound, not asserted
to be optimal.  More generally, the proof gives the lower bound
$c\min\{\Phi(\bar P),1/(Nm^2d_N^3)\}$ at every even $m\ge2d_N$,
where $d_N=\binom N2$.
\end{theorem}
\begin{proof}
Let $C$ hold with probability $1/2$ and otherwise conjugate every gate
by a uniform transposition.  On each relabeling orbit it is a quotient
of the lazy random-transposition walk on $S_N$, so its gap is at least
$c/N$ by the same character-ratio bound.  We compare its cut flow to $P$.

Every word of length $m\ge2d_N$ has at least
\[
 M(w)=\sum_{s\in T_N}\binom{n_s(w)}2
 \ge\frac{m(m-d_N)}{2d_N}\ge\frac{m^2}{4d_N}        \tag{QD.3}
\]
pairs of equal gates, by Cauchy--Schwarz.  Choose one such pair of
positions $i<j$ uniformly.  Forward Hurwitz moves bring the selected
gates, unchanged, to positions one and two, producing $(s,s,U)$.
This fixed braid is bijective and equivariant under relabeling.

To conjugate by $h\in T_N$, refresh $(s,s)$ to $(h,h)$.
Move the second $h$ right through $U$ using
$(h,u)\mapsto(huh,h)$, obtaining $(h,hUh,h)$.
Next move the last gate left through $hUh$ using inverse Hurwitz moves
$(v,a)\mapsto(vav,v)$.  The traveling gate returns to $h$, because
the product of $hUh$ is identity.  The resulting word is $(h,h,hUh)$.
Refresh the front square to $(hsh,hsh)$ and undo the original extraction
braid.  Equivariance makes the endpoint exactly $hwh$.
The route has at most $8m$ moves, all allowed by $P$.

For completeness, a deliberately loose congestion bound suffices.
Fix the positions $i,j$, the original repeated gate $s$, the conjugator
$h$, and a stage of the route.  The current word recovers the original
word uniquely: Hurwitz stages are bijections, and the only erased gate
value, $s$, is retained in this encoding.  There are at most
$4m^3d_N^2$ such encodings.  Each path has stationary weight at most
$\pi(w)/(2d_N\min M)$, while every used directed $P$ edge has capacity
at least $\pi(w)/[2(m-1)d_N]$.  Uniformity and (QD.3) give
\[
 Q_C(A,A^c)\le\rho Q_P(A,A^c),\qquad
 \rho\le16m^2d_N^3.                                \tag{QD.4}
\]
This is cut congestion, so no further path-length factor is needed.

Write $a(O)=\pi(A\mid O)$ on each relabeling orbit and
$D=\sum_O\bar\pi(O)a(O)(1-a(O))$.  The within-orbit gap of $C$ gives
$D\le C_1NQ_C(A,A^c)\le C_1N\rho Q_P(A,A^c)$.
Lumpability and Jensen's inequality also give
\[
 Q_P(A,A^c)\ge\frac12\sum_{O,O'}\bar\pi(O)\bar P(O,O')
                    |a(O)-a(O')|
 \ge\Phi(\bar P)\int_0^1\min\{\bar\pi(a>t),1-\bar\pi(a>t)\}\,dt.
\]
For $p=\pi(A)\le1/2$, the integral is at least $p-2D$.
Indeed it is the minimum of $\mathbb E|a-c|$ over constants $c$,
whereas $\min_c\mathbb E|1_A-c|=p$ and
$\mathbb E|1_A-a|=2D$.  If $D\le p/4$, the cut ratio is at least
$\Phi(\bar P)/2$; otherwise it is at least $1/(4C_1N\rho)$.
This proves the lower bound.  Lifting quotient cuts proves the upper
bound.  At the stated quadratic length $Nm^2d_N^3=O(N^{11})$.
\end{proof}

\begin{proposition}[A polynomial slowdown survives relabeling]
\label{prop:quadratic-slow-mode}
Both the full walk and its relabeling quotient satisfy
\[
 \operatorname{gap}\le\frac{C}{mN^2},\qquad
 t_{\rm mix}(1/4)\ge c mN^2=\Omega(N^4).            \tag{QD.5}
\]
\end{proposition}
\begin{proof}
Let $f(\sigma)=\#\operatorname{Fix}(\sigma)-1$ and choose $N$ consecutive
central times.  Set $F(w)=\sum_j f(G_j)$ over these times.  This observable
is invariant under simultaneous relabeling.  Compare the central segment
to an unconditioned transposition walk starting from a parity-uniform
permutation.  Its density under the identity bridge is bounded above
and below by absolute constants, by (CB.7) on the two long exterior
segments.  Under the reference law, $\mathbb E f=0$,
$\mathbb E f^2=1$ (for $N\ge5$), and
\[
 \mathbb E[f(G_{j+r})\mid G_j]
       =\left(1-\frac2{N-1}\right)^r f(G_j).
\]
The last identity follows by counting how a transposition creates or
destroys fixed points.  Summing covariances over $N$ times gives
$\operatorname{Var}F=\Omega(N^2)$.  A bounded density comparison
preserves this lower bound, using $\operatorname{Var}F=\min_c\mathbb E(F-c)^2$.

A pair update changes at most one prefix, and can affect $F$ only at
one of its $N$ selected bonds.  Given such a bond, its prefix changes
on at most four inputs.  Under the comparable parity-uniform reference
law, an affected input is fixed before or after the move with
probability $O(1/N)$, and the squared change is bounded by $16$.
Hence $\mathcal E_P(F,F)=O((N/m)/N)=O(1/m)$.
The variational bound proves (QD.5), also on the quotient because $F$
is invariant.  The standard lazy-chain relaxation lower bound gives
the mixing statement.
\end{proof}

The remaining quadratic-length target is therefore an ordered,
relabeling-invariant feature beyond a collection of well-separated
prefix permutations.  The preceding lower bound is polynomial, so it
does not by itself establish superpolynomial slow mixing.  The
base-length reduction in Theorem~\ref{thm:base-length-reduction} gives
a second way to study this same unresolved regime through arbitrary
endpoints at length $\Theta(N\log N)$.
